% ============================================================
%  Niels Treschow, "Almindelig Logik"
%  Kiøbenhavn: Trykt paa Fr. Brummers Forlag, hos Ernst Andreas Herman
%    Møller, Hof- og Universitetsbogtrykker, 1813.
%  TRANSCRIPTION (Danish) — SELECTION: §§ 26–34, printed pp. 91–105.
%
%  WHY THIS SELECTION. The 1810 essay "Gives der noget Begreb ... om enslige
%  Ting?" asserts that there is a fixed concept of singular things but never
%  says how such concepts are formed or told apart from spurious ones. These
%  §§ are where the Logik supplies the machinery: the doctrine of description
%  (§28), the higher/lower concept hierarchy with extension and comprehension
%  (§29), the distinction between PROPER and IMPROPER genera and species
%  (§30) — criteria "væsentlige og bestandige" vs "blot tilfældige og
%  ubestandige", i.e. a naturalness criterion for classifications — whole and
%  parts (§31), and the theory of division with its necessary rules (§§32–34).
%  See ../enslige-ting/ for the 1810 essay, transcribed and translated.
%
%  Scan: ~/bibliotek/Treschow, Niels/almindelig-logik-1813.pdf
%        (345 scan pp.; built from nb.no IIIF, URN:NBN:no-nb_digibok_2008040210002;
%         public domain, "Tilgang for alle").
%  OFFSET (VERIFIED in this region): scan page = printed + 8.
%        §26 opens printed 91 = scan 99; §34 runs to printed 105 = scan 113,
%        where §35 "Andre Regler" begins (not included).
%        NB the offset DRIFTS elsewhere in the volume (+10 near the end);
%        re-verify before working outside pp. 91–105.
%  TYPE: FRAKTUR. Emphasis = letterspacing (Sperrsatz) -> \emph{}.
%        Section numbers and headings are centred; headings letterspaced.
%  ORTHOGRAPHY: this book uses "ø" (NOT the "ö" of the 1810 antiqua essay).
%        Preserve: Kiendemærker, Gienstand, giøre, Slægt, Beskaffenheder,
%        nøiagtig, Spørgsmaal, Oplysning, enslige, Afdeling.
%  SECTION MAP (printed pages):
%     26. Ordforklaringer ................................. 91
%     27. Hvilke Ord man bør forklare ..................... 93
%     28. Beskrivelser .................................... 94
%     4de Afdeling. Om Delinger ........................... 95
%     29. Høiere og lavere Begreber ....................... 95
%     30. Egentlige og uegentlige Slægter eller Arter ..... 97
%     31. Et Heelt og dets Dele ........................... 98
%     32. Logiske og andre Delinger ...................... 100
%     33. Nytten af Delinger ............................. 102
%     34. Overalt nødvendige Regler ...................... 104
%  See RESUME-NOTES.md for progress and the working method.
% ============================================================
\documentclass[12pt,a4paper]{article}
\usepackage[utf8]{inputenc}
\usepackage[T1]{fontenc}
\usepackage{libertinus}
\usepackage{libertinust1math}
\usepackage[danish]{babel}
\usepackage[protrusion=true,expansion=true]{microtype}
\usepackage{setspace}
\usepackage[margin=1.2in]{geometry}
\usepackage{fancyhdr}
\usepackage{hyperref}

\hypersetup{
  pdftitle={Almindelig Logik (1813) — §§ 26–34},
  pdfauthor={Niels Treschow},
  hidelinks
}

\pagestyle{fancy}
\fancyhf{}
\rhead{Treschow, \emph{Almindelig Logik} (1813), §§ 26--34}
\cfoot{\thepage}

\setstretch{1.3}

% Centred section number + letterspaced heading, as in the print
\newcommand{\logsec}[2]{%
  \bigskip\begin{center}{\large #1}\\[4pt]{\emph{#2}}\end{center}\smallskip}

\title{\textbf{Almindelig Logik}\\[6pt]
  \large §§ 26--34 (trykte s.\ 91--105)\\[4pt]
  \normalsize Slutningen af Afsnittet om Definitioner, og\\
  \large 4de Afdeling: Om Delinger}
\author{Niels Treschow}
\date{\small Kiøbenhavn: Fr.\ Brummers Forlag, 1813}

\begin{document}
\maketitle
\thispagestyle{fancy}

\bigskip

% === TEXT BEGINS (printed p. 91 = scan 99), mid-sentence in § 25 ===

% printed p. 91
\noindent
[\dots] der, undertiden ogsaa dets Nødvendighed, som, at man ved Substans maa
forstaae den sidste indvortes og reale Grund til alle Beskaffenheder i en
Gienstand: \emph{deels at de anførte Kiendemærker ei alene forekomme, men
tillige i alle de Tilfælde, der betegnes med et vist Ord. Det sidste bevises
igien enten ved Induction af flere Exempler eller ved Deduction,} naar man af
et almindelig tilstaaet Kiendemærke udleder det, man i en real Definition helst
vil lægge til Grund. For Ex.\ deraf, at Dyd efter Talebrugen hverken tillægges
Gud, Engle eller Helgene, kan sluttes, at Dyd er moralsk Fuldkommenhed i Kamp
med Sandselighed og en stridig Natur.

\logsec{26.}{Ordforklaringer.}

\noindent
\emph{Ordforklaringer} --- ei det Samme som nominale Definitioner § 22 --- ere
saadanne, der enten af andre enstydige Ord eller af Ordets Sammensetning og
Oprindelse oplyse dets Bemærkelse, f.~Ex.\ Philosophiens, Ontologiens.
\emph{Hensigten deraf er deels at bevise dennes Overeensstemmelse med Ordets
oprindelige og sædvanlige Betydning.}
% sic: kun eet "deels"; det følgende "Det Sidste" forudsætter et andet led
Det Sidste gaaer især an med afledede og sammensatte. Thi, da de, som først
have dannet dem, formodentlig deri rettede sig efter det naturlige Forvant%
% printed p. 92
skab, de syntes at opdage mellem Stammeordenes Bemærkelse og det ny Begreb, de
vilde udtrykke; saa kan man af hine gierne slutte til hvad de ved de
sammensatte og afledede have tænkt. Etymologien tiener altsaa til at udfinde
hvad Ordene egentlig have betydet i det ældre Sprog, og efter deres Oprindelse
kan betyde. Stammeordene kan man derimod ei anderledes forklare end, at man ved
Sammenligning med de afledede viser hvad hine rimeligviis have bemærket, hvis
det nu kan være tvivlsomt eller glemt. Thi, at enten det hele Ord eller hver
Bogstav og Stavelse skulde have sin egen Betydning, det er neppe muligt at
bevise. Desuden tiener dette Slags Forklaringer til at vise den menneskelige
Forstands egen Fremgang i visse Begrebers Uddannelse, f.~Ex.\ om Skiebne, Gud og
aandelige Væsener; hvilket undertiden er meget lærerigt. Af mindre Vigtighed er
vel saadanne Ords Etymologie, der ingen videnskabelige Gienstande betegne, men
ikkun saadanne, som i det daglige og borgerlige Liv forekomme, saasom Templer,
Redskaber, Øvrigheds Personer o.~s.~v. Alligevel tiene de baade til at forklare
mange Steder hos Skribenterne, og giøre os nærmere bekiendte med Oldtidens Aand
og Skikke, som det kan være Umagen værd at kiende.

% printed p. 93
\logsec{27.}{Hvilke Ord man bør forklare.}

\noindent
Hvor Definitioner eller Ordforklaringer bør anbringes maa Enhver efter
Foredragets Øiemed selv vide at bedømme. \emph{I Almindelighed bør alle
Hovedbegreber,} d.~e.\ de, paa hvilke Læresetninger og Regler siden blive
grundede, nøiagtig forklares; da den mindste Forsømmelse deri ofte leder til
urigtige Slutninger. Enhver Definition indeholder nemlig en Setning, der lægges
til Grund for alle de Betragtninger, man over Gienstanden kan anstille. Det er
derfor nødvendigt, at denne Grundvold bliver vel lagt, om ei alt hvad man
derpaa bygger skal falde. Utydelige Begreber, skiøndt de for Resten kan være
baade klare og rigtige, f.~Ex.\ om Dyd, Pligt og Ret, sette os dog ei istand til
at besvare hvert Spørgsmaal, hvis Opløsning det selskabelige Liv og Enhvers
egen Sindsroe kan forlange, uagtet de, der studere Lovkyndigheden, i
Begyndelsen tidt ansee Definitioner paa saa bekiendte Ting for overflødige, og
fast Enhver indbilder sig forud at vide hvad han ved disse Ord bør forstaae.
Ligeledes er det med Tid, Rum og andre Ting beskaffet, om hvis sande Væsen
Philosopherne tvivle, medens intet synes Uvidende lettere at forklare.

\emph{Paa den anden Side kan denne Flid ogsaa overdrives,} naar man til ingen
Nytte op%
% printed p. 94
fylder Lærebøgerne med Definitioner over sette Begreber eller overflødige
Kunstord, men ofte løselig behandler de vigtigste, ja vel endog gaaer dem med
Taushed gandske forbi. Ingen Videnskab er bleven mere besværet med dette Slags
Overflod end Metaphysiken eller Logiken selv; hvorfra den og er indført i andre
Videnskaber. Den gode Smag har søgt at rense dem derfra. Kun Skade! at den
tilligemed saadan Ordkram undertiden har udfeiet adskilligt Godt og Grundigt,
som den mellem Skarnet ei er bleven vaer.

\logsec{28.}{Beskrivelser.}

\noindent
\emph{Beskrivelser} ere egentlig bestemte til Forklaringer over enslige
Gienstande, hvis væsentlige og bestandige Kiendemærker, som hos Enslige, enten
mangle eller ere aldeles ubekiendte. Man kan derfor ingen andre end uvæsentlige
og foranderlige anføre. Intet uorganisk Legeme synes at have sand
Individualitet; da ingen Deel deri er bestandig, og Formen selv paa utallige
Maader kan afvexle, uden at man kan fastsette nogen Punct, hvor det ophører at
være det samme, f.~Ex.\ Noahs Ark eller Theseus's Skib. I Steder og Bygninger
ere alle Kiendemærker, indtil Beliggenheden selv, som hos Tyrus, ubestandige.
Mennesker og de fuldkomnere Dyr i det mindste ere uden Tvivl virkelige
Individuer, som under alle Livets Afvexlinger udvortes og
% printed p. 95
indvortes beholde en vis uudslettelig Charakter; men hvem er istand til at
bestemme den? Vi nødes derfor til blot at beskrive dem efter Navn, Stand,
Alder, Bedrifter, Herkomst o.~s.~v., uagtet Alt dette kan forgaae, men det
enslige Menneske eller Dyr endda blive det samme. Man pleier vel, især hvad
Mennesker angaaer, at føie nogle Charaktertræk til, men disse maae altid blive
deels uvisse, deels ubestemte, de væsentlige ere lettere ved utydelige Vink at
give Leilighed til at ahne end nøiagtig at forklare.

Ved almindelige Begrebers eller Slægters og Arters Beskrivelse anføre
Naturforskerne, foruden det væsentlige eller kunstige Kiendemærke, ogsaa alle
Slags foranderlige og tilfældige. Thi da de tilfældige Beskaffenheder dog
tildeels have deres Grund i noget Bestandigt og Væsentligt; saa kan hine endog
tiene til en grundigere Kundskab om Tingenes Væsen selv. Desuden falde de
snarest i Øinene, og ere derfor lettere at udfinde.

\bigskip
\begin{center}
  {\Large 4de Afdeling.}\\[6pt]
  {\large \emph{Om Delinger.}}
\end{center}

\logsec{29.}{Høiere og lavere Begreber.}

\noindent
Et høiere Begreb er et almindeligt, hvorunder andre, der kaldes lavere, kan
tænkes indbe%
% printed p. 96
fattede, som Mennesker under Dyr, Triangler under Figur. Et Begreb, der kun
indeholder saadanne Kiendemærker, hvori flere Ting ligne hverandre, men selv
indbefattes under et høiere, kaldes i Henseende til samme \emph{Art}, det
høiere derimod \emph{Slægt}. Det samme almindelige Begreb kan altsaa i
forskiellig Henseende betragtes baade som Slægt og Art. Den høieste Slægt og
laveste Art giøre alene derfra en Undtagelse; da hiin omfatter alle, uden selv
at staae under noget høiere Begreb, og følgelig aldrig kan betragtes som Art,
denne omfatter derimod kun enslige Ting, og kan derfor ei betragtes som Slægt.

\emph{De høiere Begreber tillægger man større Omfang} --- Extension --- eller
Strækning, fordi de indbefatte flere Enslige; \emph{de lavere mere Indhold eller
Stof} --- Comprehension --- , da de, foruden de fælles Kiendemærker, have mange
særegne.

\emph{Begreber ere hverandre enten underordnede eller samordnede.} Har det
første Sted, som mellem Slægter og Arter; saa kaldes hine underordnende, disse
egentlig underordnede: hvis derimod det andet, saa ere de enten begge
indbefattede i samme Slægter, som nærmeste Arter, f.~Ex.\ Marsvioler og
Stedmodersblomstre, eller høre til forskiellige Slægter, f.~Ex.\ de nysanførte
og Natviolerne, d.~e.\ disparate. \emph{Arterne kan være deres Slægter saavel
middelbare som umiddelbare underordnede:} i sidste Til%
% printed p. 97
fælde kaldes disse deres \emph{nærmeste}, i første deres \emph{fiernere
Slægter}. Philosophien tillader sig heri mange Spring, f.~Ex.\ fra Menneske til
Dyr, der i Naturbeskrivelsen vilde være utilgivelige.

Den høieste Slægt og laveste Art er enten den \emph{absolut høieste og laveste},
saasom Ting eller Væsen, Menneske, Hest, Hund o.~s.~v.\ eller \emph{relativ} i
denne eller hiin Videnskab. Saa er Legeme i Naturlæren den høieste Slægt, i
Metaphysiken en meget lavere. I Anthropologien, hvor der handles om de
forskiellige Menneskeracer, Nationer, Kiøn, Aldre o.~s.~v.\ er Begrebet om et
Menneske det høieste, som i den almindelige Philosophie er det allerlaveste.

\logsec{30.}{Egentlige og uegentlige Slægter eller Arter.}

\noindent
\emph{Slægter saavelsom Arter ere deels egentlige, deels uegentlige eller
analogiske.} Hines Kiendemærker ere væsentlige og bestandige, disses blot
tilfældige og ubestandige. \emph{Endskiøndt man baade i Philosophien og andre
Videnskaber helst søger at opfinde det første Slags}, fordi Tingenes væsentlige
Egenskaber ere os mest magtpaaliggende at kiende; \emph{saa er dog ogsaa det
sidste Slags nødvendige}, naar der enten
% printed p. 98
ingen anden gives § 21 og 23, eller disse i andre Henseender, fornemmelig i
almindelige Theoriers Anvendelse, kan være til Nytte. Saaledes er det til
Klogskab i Omgang med Mennesker ei nok at kiende den menneskelige Natur i
Almindelighed: de forskiellige Tilbøieligheder og Omstændigheder, hvoraf de
lade sig beherske, maae ligeledes være os bekiendte. Nu er det alligevel
umuligt at udstudere hver enslig Persons Gemytsbeskaffenhed eller Charakter: og
følgelig have de Lærde giort os en Tieneste, som efter Temperamenter, Kiøn,
Alder, Herkomst o.~s.~v.\ have deelt dem i flere uegentlige Arter, hvorved de
vigtigste Forskielligheder, som hos enslige Personer kan have Sted, dog maae
blive bemærkede. Af samme Aarsag blive i Naturbeskrivelsen ei alene alle de
Arter anførte, der ansees for bestandige, men endog mange Afarter, f.~Ex.\ af
Madurter, Haveblomstre og andre Naturproducter, der baade i theoretisk
Henseende og endnu mere for Brugens Skyld kan være mærkværdige. Mange høiere
Slægter, Klasser og Ordener selv ere der uegentlige: hvoraf de saa kaldte
\emph{kunstige Systemer}, modsatte de naturlige, have deres Oprindelse.

\logsec{31.}{Et Heelt og dets Dele.}

\noindent
\emph{Et Heelt er en Enhed, der med flere Ting, som kaldes Delene,
tilsammentagne udgiør det Samme:}
% printed p. 99
f.~Ex.\ Verden med alle Himmellegemer og hvad disse indeholde, Staten med alle
dens Borgere, Slægten med alle de Arter, som dertil høre. Disse Dele kan nu
enten ligge saaledes i og uden for hverandre, at de formedelst en virkelig
eller real Forbindelse i Gienstanden selv maae betragtes som uadskillelige:
f.~Ex.\ i Verden og alle, især organiske Legemer, hvori ei alene hver Deel har
nogen Indflydelse paa de øvrige, men ingen engang ved sig selv kan bestaae;
eller dette Forhold kan alene være til i vore Forestillinger. Saaledes er der
ei mellem alle Dyr af en Art just noget Samfund, som giør det ene afhængigt af
det andet. Alligevel udgiør enhver Art et Heelt. I første Tilfælde ere de
\emph{reale}, i det andet \emph{ideale Dele}. Enslige Væsener ere Arternes,
disse Slægternes, alle tilsammentagne Dele af den høiere Slægt, men blot i
Tankerne. \emph{Denne Forskiel er meget vigtig.} Man har ofte blandet den reale
Totalitet, som er Verden, med den ideale, som er Ting, Væsen i Almindelighed,
og begge igien med den høieste Enhed i Principier eller den sidste Aarsag til
alle Ting, hvoraf de som Virkninger, hænge. Pantheismen og Emanations Systemet
beroe fornemmelig paa denne Misforstaaelse. Det høieste Væsen er hverken et
almindeligt Begreb, hvorunder Aand og Legeme, som Arter, indbefattes, ei heller
noget Heelt, som Verden, hvoraf alle Individuer ere Dele. Men, for at udrydde
saadanne Forestillinger, er det nødvendigt

\end{document}
